\chapter{Introduction}

AI systems such as chatbots and agents built on large language models (LLMs) are
now used almost everywhere. They answer customer questions, help people write
code, power search, and automate everyday tasks. As these systems spread, so do
the ways they can be misused. An attacker can trick a model into producing unsafe
or confidential output by using carefully written prompts. Two common examples are
\emph{jailbreak} prompts, which make a model ignore its safety rules, and
\emph{prompt injection}, where hidden instructions take control of the model's
behaviour.

Keeping these systems safe is hard because the threats keep changing. A one-time
safety check is not enough. New attack tricks appear all the time, and a model
that was safe last month may become unsafe after an update.

The usual approach to this problem is fragmented. Red teams (attackers) test for
weaknesses only now and then. Blue teams (defenders) respond after an alert has
been raised. The lessons learned from each test are not reliably turned into
stronger defense. The two sides rarely work together inside one system that keeps
improving on its own.

\textbf{PurpleAgents} is proposed to fix this. It is an autonomous
\emph{purple-team} framework that continuously attacks and defends a target AI
model inside a safe, authorized environment. It checks every conclusion using
two independent lines of reasoning, and it learns from every missed attack so
that the defense keeps getting better.

\section{Motivation}

Three trends make this project timely.

\begin{itemize}
  \item \textbf{AI is now a front door.} Many organisations place a chatbot or an
  AI agent directly in front of their customers, staff, or internal data. A single
  successful jailbreak can expose private information, produce harmful content
  under the organisation's name, or make an agent take an action it should not.
  The cost of a missed attack is therefore real and immediate.

  \item \textbf{Attacks are getting cheaper.} Recent research shows that LLM-based
  agents can plan and carry out attacks on their own~\cite{fang2024hack,
  xu2025forewarned}. When attacks can be automated, defense that depends on
  occasional human testing falls behind. Defense must also become continuous and
  automatic.

  \item \textbf{Single signals are unreliable.} A safety filter or an LLM ``judge''
  can be wrong. If a system trusts one signal, it can miss an attack or block a
  harmless request without anyone noticing. Checking every decision in two
  independent ways is a simple, practical way to catch such mistakes.
\end{itemize}

The literature already describes the idea of a purple agent that passes lessons
between attacker and defender~\cite{li2025shields}, but it does not build one with
verified reasoning and a learning loop. This gap, together with the trends above,
is the motivation for PurpleAgents.

\section{Existing System}

Today's AI-safety and testing tools each cover only one part of the problem.

\begin{itemize}
  \item \textbf{Red-team / attack tools} can automatically generate jailbreaks and
  injection prompts. However, they only \emph{find} weaknesses. They do not defend
  the system, and they do not learn from what they find.
  \item \textbf{Defensive tools} such as guardrails and output filters can block
  some unsafe responses. However, they react to single signals and are often
  bypassed when an attacker uses new phrasing.
  \item \textbf{Manual purple teaming} brings humans into the loop. However, it is
  slow, happens only periodically, and depends heavily on expert effort.
\end{itemize}

In summary, existing systems suffer from the following problems:

\begin{itemize}
  \item \textbf{Fragmentation:} attack and defense are kept separate, so lessons
  are not shared between them.
  \item \textbf{Single-path decisions:} one alert or one score is trusted, so a
  wrong call or a missed attack goes unchecked.
  \item \textbf{No structured learning:} a missed attack is logged as a failure,
  not turned into a reusable lesson.
  \item \textbf{Periodic, manual effort:} testing happens occasionally, not
  continuously.
\end{itemize}

\section{Problem Statement}

There is a need for an intelligent system that can continuously attack an AI
model, defend it, verify its own conclusions, and learn from every missed attack
--- all on its own, in a safe and authorized environment. Current tools cannot do
this. They solve only one piece of the problem (attack, or defense, or
reasoning), they rely on a single reasoning path, and they do not convert
failures into verified, reusable lessons. This gap leaves AI systems exposed to
new and evolving attacks.

\section{Objectives}

The objectives of this project are:

\begin{enumerate}
  \item To build an autonomous purple-team framework in which attack and defense
  agents run continuously against a target AI model.
  \item To introduce a \textbf{dual-reasoning engine}, so that a conclusion is
  accepted only when an evidence-first path and a hypothesis-first path agree.
  \item To detect attacks and, for every missed attack, find the exact reason it
  was missed.
  \item To generate testable defensive improvements, such as stronger guardrail
  rules.
  \item To keep an auditable memory of evidence, reasoning, confidence, and
  outcome.
  \item To close the loop, so that each missed attack becomes a verified lesson
  that strengthens the defense, with human approval on key changes.
  \item To show, through repeated attack--defend cycles, that the defense
  measurably improves over time.
\end{enumerate}

\section{Applications}

PurpleAgents is intended for any organisation that builds, deploys, or protects
AI systems. The main users are:

\begin{itemize}
  \item \textbf{Companies deploying chatbots and AI assistants.} Banks, online
  stores, hospitals, and service providers that put a chatbot in front of their
  customers can use PurpleAgents to test it continuously before and after each
  update, and to strengthen its guardrails over time.

  \item \textbf{AI and LLM providers.} Teams that build or fine-tune language
  models can use the framework as an automated safety check during development. A
  model can be attacked in repeated rounds, and the lessons can be turned into
  better safety rules before release.

  \item \textbf{Security Operations Centre (SOC) teams.} SOC analysts who are
  responsible for AI systems can use the auditable memory as a clear record of
  what was attacked, how it was judged, and what was fixed. Low-confidence cases
  are sent to them for review instead of being decided silently.

  \item \textbf{Developers of AI agents.} Agents that read web pages, e-mails,
  or documents are exposed to prompt injection. PurpleAgents can test such agents
  against injection attacks inside a safe environment.

  \item \textbf{Education and research.} The framework can be used in teaching
  and research to show how attack, defense, and reasoning work together, and to
  compare different guardrail strategies in a fair, repeatable way.
\end{itemize}

\section{Organization of the Report}

The rest of this report is organized as follows. Chapter~\ref{ch:lit} reviews
related research on autonomous attackers, autonomous defenders, and reasoning
methods. Chapter~\ref{ch:problem} defines the problem and the scope of the
project. Chapter~\ref{ch:req} lists the functional, non-functional, hardware, and
software requirements. Chapter~\ref{ch:design} presents the system design, the
proposed architecture, the dual-reasoning engine, and the algorithm.
Chapter~\ref{ch:analysis} covers the tools, the data sources, and the evaluation
plan. Chapter~\ref{ch:feas} presents the feasibility study and the project timeline,
and Chapter~\ref{ch:concl} states the expected outcomes and concludes the report.
